\section{Experiments}
\label{sec:experiments}


\parsection{Implementation.}
We implement our method in PyTorch~\citep{paszke2019pytorch} and train on 128 H100 GPUs with $V \in [2, 18]$ views per scene at 504-pixel longest-edge resolution. Each training step uses up to $4{,}608$ images and a fixed tokens budget ($\approx$2.5M tokens).
We use a DINOv2 ViT-B encoder with embedding dimension $768$, patch size $P{=}14$, and $R{=}4$ register tokens. The ray and depth decoders each use embedding dimension $384$ and two transformer blocks.
We sample $K \sim \text{Beta}(2, 1)$ scaled into $[8, 16]$ during training so that a single checkpoint supports any step count in this range (\cref{tab:diag_iterative}). We run $K_\text{inf} = 16$ steps at inference.
We optimize with AdamW~\citep{loshchilov2018decoupled} at a base learning rate of $3 \times 10^{-4}$, weight decay $0.05$, and a cosine decay schedule without warmup in the first stage, applying a $0.1\times$ multiplier to the DINOv2 backbone.
The first stage trains end-to-end for 200K iterations. The second stage finetunes the depth decoder (see \cref{sec:method_training}) for 40K iterations at $1 \times 10^{-4}$ with a 500-step linear warmup and a confidence regularizer weight of $\lambda_c = 0.2$.
The depth, ray, pointmap, and camera losses are weighted equally, and within the camera loss, the translation, rotation, and field-of-view terms are weighted $(1, 1, 0.5)$. 
We apply the depth-gradient term only on synthetic data, where ground-truth depth is dense enough for reliable gradient supervision. We train on a mixture of 29 public datasets and list them in the supplemental material.

\subsection{Comparison with State of the Art}
\label{sec:sota}

\parsection{Evaluation datasets.}
We evaluate on a set of diverse benchmarks that span indoor, outdoor, object-centric, and driving scenes:
DTU~\citep{jensen2014dtu},
ETH3D~\citep{schops2017eth3d},
7-Scenes~\citep{shotton2013scene},
ScanNet++~\citep{yeshwanth2023scannetpp},
and nuScenes~\citep{caesar2020nuscenes}. For ScanNet++, which serves as a training dataset for multiple baselines~\citep{lin2025depth, wang2025pi, keetha2026mapanything, duisterhof2024mast3rsfm, leroy2025grounding, vggt_omega} and our model, we enforce a clean scene-level split between the training and evaluation data.

\parsection{Evaluation metrics.}
We evaluate reconstruction quality with two metrics computed on the global pointmap after a Sim(3) alignment of the predicted points to the ground truth. Both are derived from the per-point relative error $r_i = \lVert \mathbf{X}_{\theta,i} - \mathbf{X}_{\text{gt},i}\rVert / \lVert \mathbf{X}_{\text{gt},i}\rVert$: the relative $\ell_2$ distance (\textbf{Rel.~L2}~$\downarrow$) is its mean over valid points, and the inlier ratio (\textbf{IR}~$\uparrow$) is the fraction of points with $r_i < 3\%$.
For camera pose accuracy, we report the area under the cumulative error curve at angular thresholds of $3^{\circ}$ and $30^{\circ}$ (\textbf{AUC@3}~$\uparrow$ and \textbf{AUC@30}~$\uparrow$), where the per-pair error is the maximum of the rotation and translation angle errors.

\parsection{Baselines.}
We compare against state-of-the-art feed-forward 3D reconstruction methods: VGGT~\citep{vggt}, Pi3~\citep{wang2025pi}, MapAnything~\citep{keetha2026mapanything}, and Depth Anything 3 (DA3)~\citep{lin2025depth}. We also include MASt3R~\citep{leroy2025grounding} and MASt3R-SfM~\citep{duisterhof2024mast3rsfm}, which combine pairwise prediction with sparse global alignment. In addition, we report numbers for the concurrent work VGGT-$\Omega$~\citep{vggt_omega} for completeness. We run all baselines through our evaluation framework using their official released checkpoints, with MapAnything and DA3 at their v1.1 releases, VGGT-$\Omega$ at the 1B-512 release, and DA3 reported at two backbone scales (ViT-L and ViT-G). Per-baseline configurations are detailed in \cref{sec:suppl_baselines}.





\begin{table}[t]
    \caption{
        \textbf{Model efficiency and quality.}
        Parameter count, forward-pass FLOPs (total and per-image), peak inference GPU memory, and average IR / AUC@$30^{\circ}$ across the five benchmarks of \cref{tab:comparison_sota_pmap,tab:comparison_sota_pose}, measured on a single A100 with $24$ input views.
        For MASt3R~\citep{leroy2025grounding} (\texttt{swin-5}, 120 pairs) and MASt3R-SfM~\citep{duisterhof2024mast3rsfm} (retrieval-20-10, 273 pairs), the reported FLOPs cover only the pair-network forward passes, excluding the iterative global alignment that follows. Their lower peak memory comes from processing pairs sequentially.
        \methodname{} leads on average IR and AUC@$30^{\circ}$ at the smallest parameter count.
    }
    \label{tab:comparison_efficiency}
    \footnotesize
    \setlength{\tabcolsep}{5pt}

    \centering
    \noindent\makebox[0.85\linewidth][c]{%
    \begin{tabular}{l rrrr|rr}
        \toprule
        \multirow{2}{*}{\textbf{Method}}
            & \textbf{Params} & \textbf{Compute} & \textbf{Compute / Img} & \textbf{Peak Mem}
            & \textbf{IR} & \textbf{AUC@30} \\
        & (M)~$\downarrow$ & (TFLOPs)~$\downarrow$ & (TFLOPs)~$\downarrow$ & (GiB)~$\downarrow$
        & (\%)~$\uparrow$ & (\%)~$\uparrow$ \\
        \midrule
        MASt3R~\citep{leroy2025grounding}                       & \cellcolor{tabyellow}689 & 500.0 & 20.8 & \cellcolor{taborange}4.4 & 38.3 & 63.3 \\
        MASt3R-SfM~\citep{duisterhof2024mast3rsfm}              & 690 & 1150.1 & 47.9 & \cellcolor{tabred}3.4 & 60.8 & 82.5 \\
        MapAnything~\citep{keetha2026mapanything}               & 1228 & 148.4 & 6.2 & 20.1 & 69.0 & 83.3 \\
        Pi3~\citep{wang2025pi}                                  & 959 & 153.8 & 6.4 & 6.6 & \cellcolor{taborange}77.4 & \cellcolor{tabyellow}88.6 \\
        VGGT~\citep{vggt}                                       & 1257 & 190.0 & 7.9 & 14.7 & 66.0 & 84.0 \\
        VGGT-$\Omega$-1B~\citep{vggt_omega}                        & 1144 & \cellcolor{tabyellow}99.8 & \cellcolor{tabyellow}4.2 & 7.9 & \cellcolor{tabyellow}74.8 & \cellcolor{taborange}90.7 \\
        DA3-L~\citep{lin2025depth}                              & \cellcolor{taborange}356 & \cellcolor{tabred}71.4 & \cellcolor{tabred}3.0 & 7.3 & 58.3 & 79.7 \\
        DA3-G~\citep{lin2025depth}                              & 1201 & 178.7 & 7.4 & 13.0 & 66.5 & 84.7 \\
        \midrule
        \textbf{Ours}                                           & \cellcolor{tabred}\textbf{117} & \cellcolor{taborange}\textbf{75.9} & \cellcolor{taborange}\textbf{3.2} & \cellcolor{tabyellow}\textbf{4.9} & \cellcolor{tabred}\textbf{80.3} & \cellcolor{tabred}\textbf{91.8} \\
        \bottomrule
    \end{tabular}}
    \vspace{-4mm}
\end{table}



\parsection{Results.}
At $117$M parameters and $75.9$ TFLOPs of compute (\cref{tab:comparison_efficiency}), \methodname{} matches or exceeds much larger feed-forward methods while running in under $5$ GiB of peak memory at $24$ input views.
Among prior work, Pi3~\citep{wang2025pi} is the closest competitor, leading on indoor pointmap accuracy at $8{\times}$ our parameters and $2{\times}$ our compute. VGGT~\citep{vggt} retains its lead on DTU pose at $10{\times}$ our parameters, with its lead concentrated at the tightest AUC@$3^{\circ}$ threshold ($96.5$ vs $83.2$) and shrinking to within $1.0$ point at AUC@$30^{\circ}$.
The concurrent VGGT-$\Omega$~\citep{vggt_omega} edges \methodname{} on outdoor pointmap Rel.~L2 (ETH3D, nuScenes) and on 7-Scenes pose at $10{\times}$ our parameters, but trails on every other benchmark and on the average IR and AUC@$30^{\circ}$ at the bottom of \cref{tab:comparison_efficiency}; on ScanNet++ in particular, \methodname{} leads VGGT-$\Omega$ by nearly $50$ points at AUC@$3^{\circ}$ and over $10$ points at AUC@$30^{\circ}$.
MASt3R-SfM~\citep{duisterhof2024mast3rsfm} is competitive on indoor sequences but at or near the bottom of every nuScenes metric, at $15{\times}$ our forward-pass cost. Overall, our method achieves top average performance across benchmarks, the highest parameter efficiency, and remains highly competitive in compute and memory cost.

\begin{figure*}[t]
    \centering
    \includegraphics[width=\linewidth]{figures/qualitatives.pdf}
    \caption{
        \textbf{Qualitative results.}
        Predicted point clouds for three scenes, comparing \methodname{} to four feed-forward baselines with parameter counts shown above each column. \methodname{} produces denser, less noisy point clouds despite using far fewer parameters.
    }
    \label{fig:qualitative}
    \vspace{-4mm}
\end{figure*}



\subsection{Analysis}
\label{sec:analysis}

We diagnose two axes of our recurrent backbone: architectural choices in the block (\cref{tab:diag_block}), and the step count at training and inference (\cref{tab:diag_iterative}).
All variants use a ViT-B encoder and are trained for 100K iterations on the same data, and differ only along the examined axis.


\parsection{Block design.}
We progressively add the components of our recurrent design in \cref{tab:diag_block}, starting from a fully decoupled variant in which each of the 16 recurrent steps has its own block (no weight sharing, no time conditioning).
Weight sharing collapses the 16 independent blocks into a single shared block applied recurrently. The time-conditioned residual gates ($\mathbf{s}_\text{attn}, \mathbf{s}_\text{mlp}$) modulate the block's attention and MLP branches, and the state gate ($\mathbf{s}_\text{out}$) yields our full method.
Each component improves every metric monotonically. Notably, weight sharing alone already outperforms the decoupled architecture despite having $16{\times}$ fewer parameters.
\begin{table}[t]
    \caption{
        \textbf{Block design ablation.}
        We progressively add the components of our looped block.
        Weight sharing constrains the recurrence to a single looped block, the time-conditioned residual gates modulate the attention and MLP branches, and the state gate modulates the residual stream.
        We report metrics averaged across the five benchmarks of \cref{tab:comparison_sota_pmap,tab:comparison_sota_pose}.
    }
    \label{tab:diag_block}
    \footnotesize
    \setlength{\tabcolsep}{6pt}
    \centering
    \noindent\makebox[0.85\linewidth][c]{%
    \begin{tabular}{l|ccc|rrrr}
        \toprule
        \multirow{2}{*}{\textbf{Variant}}
            & \makecell{Weight \\ sharing}
            & \makecell{Residual gates \\ ($\mathbf{s}_\text{attn}, \mathbf{s}_\text{mlp}$)}
            & \makecell{State gate \\ ($\mathbf{s}_\text{out}$)}
            & \multirow{2}{*}{\textbf{Rel.~L2} $\downarrow$}
            & \multirow{2}{*}{\textbf{IR} $\uparrow$}
            & \multirow{2}{*}{\textbf{AUC@3} $\uparrow$}
            & \multirow{2}{*}{\textbf{AUC@30} $\uparrow$} \\
        \midrule
        Decoupled & \xmark & \xmark & \xmark & 0.056 & 61.1 & 23.0 & 82.0 \\
        Shared & \cmark & \xmark & \xmark & \cellcolor{tabyellow}0.045 & \cellcolor{tabyellow}66.4 & \cellcolor{tabyellow}30.2 & \cellcolor{tabyellow}84.8 \\
        Shared + residual gates & \cmark & \cmark & \xmark & \cellcolor{taborange}0.042 & \cellcolor{taborange}67.0 & \cellcolor{taborange}31.5 & \cellcolor{taborange}85.9 \\
        \midrule
        \textbf{Shared + state gate} & \cmark & \cmark & \cmark & \cellcolor{tabred}0.040 & \cellcolor{tabred}69.2 & \cellcolor{tabred}33.3 & \cellcolor{tabred}86.9 \\
        \bottomrule
    \end{tabular}}
    \vspace{-4mm}
\end{table}


\parsection{Step-count.}
We vary the recurrent step count at training and inference in \cref{tab:diag_iterative}. The fixed-$K$ baselines train and run with a constant step count. Our variable-$K$ model trains with $K \sim \text{Beta}(2, 1)$ on $[8, 16]$ and is evaluated at $K_\text{inf}{=}12$ and $K_\text{inf}{=}16$.
At $K_\text{inf}{=}16$, variable-$K$ training matches Fixed $K{=}16$ within $\sim$2\% on every metric. The same checkpoint, evaluated at $K_\text{inf}{=}12$, stays within $\sim$3\% of Fixed $K{=}12$ across all metrics.


\begin{table}[!t]
    \caption{
        \textbf{Step-count analysis.}
        Fixed-$K$ baselines train and run with a constant step count. Our variable-$K$ model trains with $K \sim \text{Beta}(2, 1)$ on $[8, 16]$ and is evaluated at $K_\text{inf}{=}12$ and $K_\text{inf}{=}16$. Metrics are averaged across the five benchmarks of \cref{tab:comparison_sota_pmap,tab:comparison_sota_pose}.
        \methodname's variable-$K$ training stays within $\sim$2\% of Fixed $K{=}16$ at $K_\text{inf}{=}16$ and $\sim$3\% of Fixed $K{=}12$ at $K_\text{inf}{=}12$, so a single checkpoint covers both inference budgets at near-zero quality cost.
    }
    \label{tab:diag_iterative}
    \footnotesize
    \setlength{\tabcolsep}{6pt}
    \centering
    \noindent\makebox[0.85\linewidth][c]{%
    \begin{tabular}{l|cc|rrrr}
        \toprule
        \multirow{2}{*}{\textbf{Variant}}
            & \multirow{2}{*}{\makecell{Training\\ $K$-sampler}}
            & \multirow{2}{*}{\makecell{$K_\text{inf}$}}
            & \multirow{2}{*}{\textbf{Rel.~L2} $\downarrow$}
            & \multirow{2}{*}{\textbf{IR} $\uparrow$}
            & \multirow{2}{*}{\textbf{AUC@3} $\uparrow$}
            & \multirow{2}{*}{\textbf{AUC@30} $\uparrow$} \\
            & & & & & & \\
        \midrule
            Fixed $K{=}12$ & fixed & 12 & 0.044 & 67.6 & 30.4 & 85.5 \\
            \textbf{Ours} (Variable, $K_\text{max}{=}16$) & Beta on $[8, 16]$ & 12 & 0.043 & 66.7 & 29.6 & 85.4 \\
        \midrule
            Fixed $K{=}16$ & fixed & 16 & 0.041 & 69.6 & 33.8 & 86.8 \\
            \textbf{Ours} (Variable, $K_\text{max}{=}16$) & Beta on $[8, 16]$ & 16 & 0.040 & 69.2 & 33.3 & 86.9 \\
        \bottomrule
    \end{tabular}}
    \vspace{-4mm}
\end{table}













